% The schema the four editions of Section 8 share, and what it becomes here.
%
% Drawn rather than generated, because it states decisions and not data. The
% upper row is the schema: five stations and one loop, the same in all four
% editions. The lower row, in grey, is the same schema with this edition's
% values filled in. The siblings fill it in differently (Table 3), which is
% the point of Section 8: the stations do not move, the values do.
\begin{figure}[htbp]\centering
\resizebox{\textwidth}{!}{%
\begin{tikzpicture}[
  font=\small,
  st/.style={draw, rounded corners=2pt, align=center, inner sep=4pt,
             minimum height=12mm, text width=27mm},
  src/.style={st, fill=black!3},
  inst/.style={font=\scriptsize\itshape, text=black!55, align=center,
               text width=27mm, anchor=north},
  lbl/.style={font=\scriptsize, text=black!55, align=center},
  arr/.style={-{Stealth[length=4pt]}, draw=black!55},
  human/.style={arr, dashed},
]

% The schema: five stations.
\node[src] (W) {\textbf{Witness}\\[1pt]\scriptsize held by the institution,\\\scriptsize never stored};
\node[st, right=17mm of W] (R) {\textbf{Reading}\\[1pt]\scriptsize the one file written,\\\scriptsize in a declared language};
\node[st, right=17mm of R] (D) {\textbf{Derivations}\\[1pt]\scriptsize computed from the\\\scriptsize reading, on every build};
\node[st, right=17mm of D] (V) {\textbf{Verifications}\\[1pt]\scriptsize what a machine can\\\scriptsize check of them};
\node[st, right=17mm of V] (A) {\textbf{Address}\\[1pt]\scriptsize the reading beside\\\scriptsize the witness, cited};

\draw[arr] (W) -- node[lbl, above=2pt]{the model} node[lbl, below=2pt]{a pass} (R);
\draw[arr] (R) -- node[lbl, above=2pt]{scripts} node[lbl, below=2pt]{rebuilt} (D);
\draw[arr] (D) -- node[lbl, above=2pt]{checkers} node[lbl, below=2pt]{or fail} (V);
\draw[arr] (V) -- node[lbl, above=2pt]{the site} node[lbl, below=2pt]{cited} (A);

% The loop: the only step no file can prove, and the only one that changes the reading.
\node[st, dashed, fill=black!3, below=10mm of D, text width=52mm, inner sep=5pt] (P)
  {\textbf{A person}\\[1pt]\scriptsize compares the reading with the witness,\\\scriptsize corrects it at source, and declares};
\draw[human] (A.south) |- (P.east);
\draw[human] (P.west) -| node[lbl, above=2pt, pos=0.3]{the correction} (R.south);

% The same schema, with this edition's values.
\node[inst, below=5mm of P.south, text width=0mm] (anchor) {};
\node[inst] at (W |- anchor) {Whitney photograph,\\ResourceSpace reference,\\fetched at view time};
\node[inst] at (R |- anchor) {\texttt{batch-NN.tex}\\twelve sheets a pass,\\14 macros};
\node[inst] at (D |- anchor) {TEI, PDF, HTML,\\accounts, formats,\\materials, days};
\node[inst] at (V |- anchor) {\texttt{hopper.odd}\\Lean 4 certificates\\CI regenerates and compares};
\node[inst] at (A |- anchor) {hopper.commutator.io\\sheet, leaf, pass\\each cited separately};
\node[lbl, anchor=east, xshift=-3mm] at (W.west |- anchor) {\itshape here:};

\end{tikzpicture}}
\caption{The schema the four editions of Section~8 share, and its values in this
one. Five stations and one loop: the model turns a witness into a reading, in
passes whose header names the model and the date; scripts derive everything
else from the reading and regenerate it on every build; checkers fail the build
on what a machine can check; the site places the reading beside the witness
and cites both by reference. The dashed loop is the person, who compares the
reading with the witness and corrects the reading, which alone changes anything
downstream. Four invariants hold at every station: the witness is never stored;
only the reading is written, everything else is derived and compared; every
assertion carries an address into the witness; and the person's verdict is
declared, not observed. The grey row gives this edition's values; the siblings'
are in Table~\ref{tab:params}.}\label{fig:pipeline}
\end{figure}
